% +------------------------------------+
% |   Generated by www.docx2latex.com  |
% |   Version: 2.0.0                   |
% +------------------------------------+

\documentclass[10pt]{article}

\usepackage{adjustbox}
\usepackage{amsmath}
\usepackage{amssymb}
\usepackage{fancyhdr}
\usepackage{float}
\usepackage[T1]{fontenc}
\usepackage{graphicx}
\usepackage[utf8]{inputenc}
\usepackage{mathtools}
\usepackage{multicol}
\usepackage{multirow}
\usepackage{wasysym}
\usepackage[backend=biber, sorting=none]{biblatex}
\addbibresource{bibliography-biblatex.bib}
\usepackage[paperheight=29.7cm,paperwidth=21.01cm,left=3.5cm,right=3.5cm,top=3.25cm,bottom=2.5cm]{geometry}


\setlength\parindent{0pt}
\renewcommand{\arraystretch}{1.3}\pagestyle{fancy}
\fancyhf{}
\lhead{\thepage}



\title{Frequent Losers in Public Procurement}

\begin{document}
\maketitle

\vspace{2\baselineskip}
\begin{center}
Darcio Genicolo-Martins\footnote{ We are very thankful to Rita Joyanovic, Volnir Pontes Junior, Mário Alexandre Reis da Silva, and the Department of Finance of Sao Paulo State (SEFAZ/SP) staff for their outstanding collaboration.}\textsuperscript{,} \footnote{ Assistant Professor at Insper, Center for Regulation and Democracy, Insper.E-mail: darciogm1@insper.edu.br.}
\end{center}


\begin{center}
Paulo Furquim de Azevedo\footnote{ Senior Fellow Research and Director of the Center for Regulation and Democracy at Insper. E-mail: pfazevedo@insper.edu.br.}
\end{center}


\vspace{1\baselineskip}
\begin{center}
February 2024
\end{center}


\begin{center}
Preliminary Version (Comments Welcome)
\end{center}


\vspace{1\baselineskip}

\begin{center}
{\Large\textbf{Abstract}}
\end{center}

The systematic loss by firms of public tenders in which they participate for long periods and several times may indicate a cartel. This paper proposes a screening method to detect cartels in public tenders by considering frequent losers. Using data on public procurement in Sao Paulo State, Brazil, from 2009 and 2019, we estimate that frequent losers are associated with 10$\%$ higher prices, 32$\%$ more participants, and 29$\%$ more bids. These results are consistent with the behavior of a cartel (higher prices) that tries to avoid detection by manipulating variables that signal competition (number of players and bids). The proposed method can address two limitations of traditional screening methods: (i) the ability to distinguish between tacit collusion and explicit collusion and (ii) the identification of a possible cartel before the conclusion of public tender processes.

\begin{center}
{\footnotesize \textit{Keywords: cartel; screening; public procurement; frequent losers.}}
\end{center}


\vspace{2\baselineskip}
\section{Introduction}

There is an intriguing phenomenon in public procurement. Some participants in bidding contests lose out systematically for long periods and several times. This fact is interesting because there is a fixed cost for companies to participate in public bids; therefore, under competitive conditions, it would be expected that these companies would win some opportunities or be excluded from this market.

An alternative hypothesis is that these firms, herein referred to as ‘frequent losers,’ are operating not autonomously but as part of a cartel in which it is possible to transfer the gains from the cartel among its members. Because of these characteristics, the presence of frequent losers has the potential to serve as a mechanism to detect a cartel in public bids. Usually, antitrust authorities use two main types of cartel detection mechanisms. The first method gained traction throughout the 1990s in the form of leniency agreements and award-winning disclosure  \cite{wils2007leniency}. The second mechanism consists of screening methods that became predominant from the mid-2000s due to the greater availability of data and computational capacity. Through observable data such as prices and quantities, screening methods seek to identify behaviors that are consistent with the cartel hypothesis (i.e., coordination among competitors) and inconsistent with the competition hypothesis among participants in a given market.

 Although these models might be insufficient to prove a cartel’s existence, they are fundamental to drive investigative efforts and may constitute evidence for judicial authorization of search and seizure operations and interception of communications.

Screening methods available in the literature, such as  \cite{green1984noncooperative},  \cite{haltiwanger1991the}  \cite{haltiwanger1991the}  \cite{abrantesmetz2006a}, have two fundamental limitations. First, these models cannot distinguish tacit collusion from cases of explicit collusion (i.e., cartels). This limitation is problematic for their practical application since tacit collusion, unlike cartels, does not constitute an antitrust offense  \cite{harrington2008detecting}.

\ \ \ \ \ Second, screening methods solely identify collusive behavior after its occurrence and therefore serve as support only for the repression of cartels but not for their prevention. In public procurement, it would be desirable for screening methods to identify collusive behavior before the bidding process occurs, through observable variables in public notices and participants’ registration, to reduce the social costs of cartels.

This paper proposes an alternative screening method that directly addresses these two limitations and does not require additional data beyond those already available to competition authorities and control bodies.

Primarily, the proposed method consists of the use of frequent losers as ‘flags’ for the ex-ante screening of cartels in public tenders. To this end, this research proposes a method for identifying frequent losers to differentiate these companies from others that, although they may lose frequently, do not exhibit sufficiently abnormal behavior. The article also presents frequent losers’ descriptive characteristics, reinforcing their properties as markers of collusion.

Finally, this paper examines the relationship between the presence of frequent losers in public bids and the price level and other indicators of competition intensity, such as the number of participants and the number of bids. The results indicate that tenders in which frequent losers participate have 10$\%$ higher prices, a 32$\%$ higher number of participants, and a 29$\%$ increase in the number of bids.

These results are consistent with cartel behavior (higher prices) and a strategy to avoid detection by increasing signs of higher competitive intensity. Note that this is precisely the expected effect of frequent losers: more participants and no competitive pressure on the winning bid.

This paper has five sections, including this introduction. Section 2 provides an overview of cartel screening methods to clarify the contributions of this research. Section 3 presents the method for identifying frequent losers and its descriptive characteristics. An analysis of the relationship between frequent losers and the results of public tenders is presented in section 4. Finally, in section 5, we present the implications of the results and final considerations.

\vspace{1\baselineskip}
\section{Models for Detecting Cartels in Public Tenders}

Cartels are a common practice and are quite costly to society. It is estimated that less than 15$\%$ of US cartels are identified by the Department of Justice  \cite{bryant1991price}. Considering that US anti-cartel policy is recognized as the most active, it is reasonable to assume that this proportion is even lower in countries with less antitrust enforcement tradition.

 Cartels are associated with higher prices to final consumers  \cite{connor2007price}, resulting in lower production than the socially optimum level and costs of maintaining the cartel and covering up its activities. For those reasons, the identification and subsequent punishment of cartels are a priority for virtually all competition authorities. Since it is a known illicit practice, it is not easy to detect cartels. There are two essential and complementary instruments to identify this anticompetitive behavior.

 The first instrument is disclosure by one of the cartel’s participants through leniency agreements or award-winning disclosure or of third-party claimants who had access to direct evidence of the agreement among the cartel’s participants  \cite{wils2007leniency}. Even if the competition authority develops a leniency program, its position is predominantly passive because third parties perform the identification of the cartel and the initial collection of evidence (Wils 2016).

 The second identification tool is predominantly active and consists of methods for analyzing companies’ behavior in different markets with the aim of identifying suspicious behavior that is consistent with the existence of a cartel and inconsistent with the hypothesis of competition between companies. These methods are known as \textit{screening methods}  \cite{harrington2008detecting}.

 Observing suspicious behavior is not enough to prove a cartel’s existence, but it is a crucial tool to guide competition authorities’ investigations in their search for direct evidence. Evidence can be obtained through search and seizure operations or by intercepting communications, which are only authorized by the judiciary when there is robust evidence of suspicious conduct. Screening methods could help to identify this evidence.

Those two detection instruments - leniency agreements and screening methods - are complementary, not substitutes  \cite{hüschelrath2014cartel}  \cite{schinkel2013balancing}. Companies’ incentives to opt for the leniency agreement by providing the competition authority with evidence of a cartel are associated with the probability of cartel detection. Proper screening methods can significantly increase the likelihood of identifying cartels.

\ \ In addition, the most stable cartels (i.e., those with a lower probability of detection) are possibly the most harmful to society and the least vulnerable to award-winning disclosure. For these reasons, authorities increasingly tend to use screening methods in their investigations  \cite{schinkel2013balancing}  \cite{imhof2018screening}. This trend is reinforced by increased data availability and processing capacity through machine learning and artificial intelligence, which give competition authorities a greater capacity to identify suspicious behaviors (Sanchez-Graells 2019).

\ \ The literature on cartel detection models is reasonably prolific. In general, statistical models use observed price and quantity data to infer firms’ patterns of behavior and reaction curves to identify whether such patterns are consistent with the cartel assumptions and inconsistent with the competition hypotheses.

Some of these models test whether the behavior of companies participating in a cartel differs from the action of nonparticipants, such as  \cite{porter1993detection} and  \cite{porter1999ohio}. In both cases, the model shows that the suspicion of cartel conduct, with defined participants and a defined period of operation, presents a correspondence with their market behavior, which has the value of proof if combined with a \textit{plus factor}  \cite{hovenkamp2005federal}). However, it is not a full screening method; its role is not to identify suspicious conduct but to test the relevance of evidence that may have been brought by a complaint. Full screening methods use markers to signal which markets and firms are engaging in suspicious behavior\footnote{ A comprehensive review of screening models is presented by Harrington (2008).}. These markers can be of three types.

\ \ \ The first deals with the relationship between firms’ prices and demand movements, considering that a cartel structure reacts to demand fluctuations differently than firms that operate in competition. Some of the most notable papers on screening methods are from this first group, such as  \cite{green1984noncooperative} and  \cite{haltiwanger1991the}. The second group is about the relationship between market shares and price variance and is based on theoretical models and the empirical regularity of greater price stability in cartelized markets, which is broadly well documented in the literature  \cite{abrantesmetz2006a}  \cite{imhof2019detecting}. Finally, the third group focuses on the relationship between firms’ prices, estimating reaction functions, and how they would be associated with competitive or collusive behavior  \cite{bajari2003deciding}.

\ \ \ \ \ In the case of screening methods for public procurement, the variables used for cartel detection are specific to those observed in public tenders, such as the bidding pattern and the existence or absence of rotation among the winners, as is the standard of cartel operation; it has become conventional to call this \textit{bid-rigging}  \cite{imhof2019detecting}).

 This pattern is essential in cases where the cartel participants do not have mechanisms for transferring the revenues derived from the cartel, so the rotation of winners becomes a necessary mechanism to avoid defections. However, this identification is not appropriate for cases where transfer mechanisms are feasible, either via corporate control or informally among cartel participants.

One of the difficulties that should be addressed by screening methods is their ability to rule out alternative hypotheses. For example, in the case of public bids, the bidding pattern may respond to cost differentials between companies, resulting in bids and ranking of winners that might appear to be the result of coordination between competitors.

An interesting proposal to circumvent this problem is offered by  \cite{kawai2019using}, who use the margin of victory, analogous to a discontinuous regression method, to compare firms with supposedly similar cost structures to rule out this alternative hypothesis. They show that the inferences of the model or the robustness of the results is a recurring concern of the new screening methods.

 There are two characteristics of the models mentioned that deserve to be highlighted because they emphasize the contributions of this paper. First, none of the cited models distinguishes tacit collusion from cases of explicit collusion, i.e., cartel  \cite{harrington2008detecting}). However, this distinction is fundamental to the application of antitrust policy because, although the cartel is considered the most serious of antitrust offenses, tacit collusion is not even an offense and, therefore, not punishable.

 The second characteristic is that the models mentioned above are a means of detecting a cartel after its occurrence. The variables used are prices, quantities, and the winners’ rotation pattern observed in a given period. Therefore, these methods are here called \textit{ex-post screening models}. In conventional markets (i.e., those that are not public procurement), this qualification is unnecessary because the repression of cartels is typically carried out a posteriori, and preventive intervention is performed through the control of structures such as mergers and acquisitions  \cite{hovenkamp2005federal}).

 In the case of public procurement, however, the cartel is often held before the bidding process, i.e., before the participants disclose their bids. Identifying the cartel a posteriori fulfills the punitive nature by intending to inhibit future illicit conduct, but it does not prevent the loss of a fraudulent tender by the cartel’s practice. These costs are exceptionally high in public bids, which require time for planning and high transaction costs to cancel contracts in execution and to conduct a new tender process.

Therefore, it would be desirable to develop \textit{ex-ante cartel screening methods} that could identify suspicious behavior before the execution of the bid. The use of frequent losers as a marker for collusion may consist of ex-ante screening, thereby distinguishing it from the other models in this regard. This is the main contribution of this research.

This marker can also be used in conjunction with other markers and variables in the ex-post analysis to strengthen the analysis, as recommended by  \cite{harrington2008detecting}). The ‘frequent losers’ marker might distinguish tacit collusion from explicit collusion, bringing more efficiency to the investigation techniques\footnote{ See section 4 for more details.}.

 Finally, it is relevant to note that cartelized companies respond strategically to the enforcement of antitrust authorities in what  \cite{schinkel2013balancing}) called a \textit{cat $\&$ mouse game}: cartelized companies will simultaneously seek to raise prices and to avoid detection by the competition authority. This situation may result in high prices coinciding with other market evidence that indicates increased competition, such as the number of participants and price dispersion.

 In addition, cartel participants will seek to alter their behavior according to the detection model used to avoid being caught. For this reason,  \cite{schinkel2013balancing} argues that the development of screening methods is a continuous process that is necessary every time cartelized companies change their behavior to avoid detection. Companies increase costs in this process, which may render the cartel unfeasible when avoiding detection is more costly than its benefits.

\vspace{1\baselineskip}
\section{Frequent Losers: Definition and Characteristics}

The collusion marker proposed by this paper requires an empirical definition of frequent losers. This section presents this definition, beginning by presenting the data used and followed by the method of identifying frequent losers.

Finally, this section presents a descriptive analysis of the main characteristics of this group of firms in the sample. These descriptive characteristics are already sufficient to distinguish the behavior of this group of bidders from what would be expected of bidders operating under competitive conditions.

\vspace{1\baselineskip}
\section{Data Source and Description}

We use administrative data on bidding-level public procurement tenders of common goods and services in the State of Sao Paulo, Brazil, from January 2009 to December 2019. All transactions took place under the electronic procurement platform named Bolsa Eletronica de Compras (BEC), which is available for all PBUs (Public Buyer Units) across the state. SEFAZ/SP (Department of Finance of Sao Paulo State) is responsible for the operational management and centralization of BEC’s bidding data.

\section{In total, 1,344 PBUs make regular purchases at BEC, including state-level executive, legislative, and judiciary bureaus in the State of Sao Paulo as well as other affiliated entities, such as some municipalities located in the State of Sao Paulo and a group of other private organizations. PBUs purchased 169,607 different types of items (goods and services), totaling 3,866,407 transactions from 19,007 distinct firms in this period.}

\section{BEC has a very detailed catalog of standardized goods and services organized at three levels of detail: group, class, and item. Data are organized by purchase offers (PO), the electronic document issued by the PBU that identifies and quantifies a set of goods and services that will be purchased. A PO is defined by a 22-character code and may contain one or more items listed, but each item has its own purchase process. Thus, the purchase of an item is uniquely identified by the item code purchase offer (POI), consisting of the combination of the PO and the purchased item codes.}

\vspace{1\baselineskip}
\section{Table 1. Descriptive Statistics: Public Tenders}

\vspace{1\baselineskip}
\section{For each POI, it is possible to observe parameters such as item quantities and reference prices and tender outcomes such as bid prices (winners and losers), the number of participating firms, the number of bids, whether the public tender was successfully, and the identification of the PBU and the auctioneer, among other variables.}

\section{Additionally, there is much information about the companies that participate in tenders. For each firm, it is possible to observe its uniquely defined firm national code (CNPJ), firm age, geocoded address, three-digit level of the National Classification of Economic Activities (CNAE), number of victories and losses in public tenders, number of bids and all bid values offered in every tender process, among others.}

\vspace{1\baselineskip}
\section{Definition of Frequent Losers: ‘Always Losers’ Outliers}

\section{To properly define ‘frequent losers,’ it is necessary to identify suppliers who participated in the bid at least once and never won - the ‘always losers’ firms. Figure 1 shows the number of losses of suppliers that failed in all their participation in public bids from 2009 to 2019. The data show a highly unequal distribution of defeats and victories among the companies.}

\section{The vast majority of firms lost only a few times, indicating that they participated in a small number of contests. However, it is possible to observe a group of companies that participated in many public tenders and lost systematically. As mentioned, the existence of firms that frequently lose tenders raises the question of why they continue to compete despite a growing number of defeats over time.}

\vspace{1\baselineskip}
\begin{figure}[H]
\centering
\includegraphics[width=11.24cm,height=8.43cm]{./images/image1.png}
\caption{Number of losses : 'Always Losers'}
\label{fig:number_losses__always_losers}
\end{figure}


\vspace{1\baselineskip}
\section{An outlier detection method, the interquartile range method (IQR), was chosen to identify the frequent loser sellers. This method separates the data into two groups according to a calculated threshold: (i) a standard group (companies that lost only a few times and stopped competing for whatever reason) and (ii) an outlier group (firms that participated in many tenders and won none).}

\section{The IQR algorithm classifies the outliers as firms whose number of losses is 1.5 times above the distance between the first and third quartiles plus the median of firms’ losses. For this paper, these outlier sellers are called frequent losers. There are 2,471 companies identified as frequent losers using this classification, indicated to the right of the red vertical line in Figure 2. At least one frequent loser participated in 73,591 tenders (approximately 3.61% of all bids).}

\vspace{2\baselineskip}
\begin{figure}[H]
\centering
\includegraphics[width=10.51cm,height=7.88cm]{./images/image2.png}
\caption{Identifying Frequent Losers}
\label{fig:identifying_frequent_losers}
\end{figure}


\vspace{1\baselineskip}
\section{It is possible to observe much better uniformity within this group. In addition, Figures 3(a) and 3(b) suggest that frequent losers tend to lose consistently over time. This fact reinforces the perennial character of the behavior of frequent losers.}

\vspace{3\baselineskip}
\begin{figure}[H]
\centering
\includegraphics[width=11.69cm,height=8.77cm]{./images/image3.png}
\caption{(a) Number of Months Competed. (b) Total Number of Losses.}
\label{fig:a_number_months_competed_b}
\end{figure}


\ \ \ \ 

\vspace{1\baselineskip}
\section{BEC Sample: Characteristics of Frequent Losers}

\section{The main objective of this section is to characterize frequent losers in the BEC sample. First, we provide a description of the modes of tenders in which the frequent losers participate. Frequent losers are almost equally involved in sealed-bid (convites) and multi-round descending auctions (pregoes), but they rarely participate in direct negotiation (dispensa de licitacao), as shown in Figure 4.}

\vspace{1\baselineskip}
\ \ \ \ \begin{figure}[H]
\includegraphics[width=11.87cm,height=8.9cm]{./images/image4.png}
\caption{Modes of Tenders}
\label{fig:modes_tenders}
\end{figure}
\ \ \ \ 

\section{Pregoes and convites are different procedures whose successive losses might have different interpretations. There is a relatively higher transaction cost to participate in pregoes since in this mode, it is necessary to actively participate on the day of the tender by submitting bids and, if required, presenting appeals and questions to the personnel in charge. In the case of convites, participants only need to submit one-shot proposals, which reduces the costs.}

\section{In terms of participation in different modes of tenders, it is interesting to observe that companies can also be separated into three distinct groups: (i) firms that only participate in convites; (ii) firms that only participate in pregoes; and (iii) firms that participate in both modes. Figure 5 shows that approximately 14% of frequent losers only participate in convites and 49.4% participate only in pregoes. Approximately 37% of frequent losers are present in both convites and pregoes.}

\vspace{1\baselineskip}
\textit{\begin{figure}[H]
\includegraphics[width=11.06cm,height=7.25cm]{./images/image5.png}
\caption{Types of Frequent Losers}
\label{fig:types_frequent_losers}
\end{figure}
}

\vspace{2\baselineskip}
\section{Analogously, the diversity of item groups to which frequent losers are linked may lead to another type of classification. A frequent loser can compete in tenders of many distinct groups of items or participate in bids of specific ones. Thus, the diversity (or lack of diversity) of purchase items of a frequent loser who chooses to participate might provide relevant additional information.}

\section{To deepen this analysis, the diversity of items in which companies participate, Shannon’s entropy coefficient, is calculated for each group of items and each frequent loser firm, defined according to the following equation:}

\vspace{2\baselineskip}
\begin{equation}
H = -\sum_{}^{}p\left(x\right)\log p\left(x\right)     (1)
\end{equation}


\vspace{2\baselineskip}
\section{is the final entropy value, and p (x) is the probability that each element would be found randomly in the universe of elements. If H is closer to 0, it means that companies bid for only one group of items or something close to that. The higher H is, the greater the diversity of item groups. An entropy of 5 means that the company bids for approximately 32 different groups of items.}

\section{There are at least two valid interpretations regarding the results obtained. First, firms can act in good faith in public tenders for different groups of items to diversify their activities. Second, it is also possible that companies are only used to simulate competition and manipulate the tender. Figure 6 shows the graph of the entropy of the item groups for all frequent losers.}

\vspace{1\baselineskip}
\begin{figure}[H]
\includegraphics[width=11.02cm,height=8.27cm]{./images/image6.png}
\caption{Entropy Coefficient}
\label{fig:entropy_coefficient}
\end{figure}
\ \ \ \ 

\vspace{1\baselineskip}
\section{It is possible to observe that half of the companies participate in tenders related to only one item, and approximately half participate in contests related to two or more item groups. The analysis is limited to companies that participate in tenders of less than 16 distinct item groups. The frequency at which every item appears reveals a high concentration of tenders in a few groups of items. Figure 7 presents the graph of the frequency of the item groups.}

\vspace{3\baselineskip}
\begin{figure}[H]
\centering
\includegraphics[width=12.55cm,height=9.41cm]{./images/image7.png}
\caption{Item Group Variety}
\label{fig:item_group_variety}
\end{figure}


\vspace{1\baselineskip}
\ \ \ \  

\section{The high concentration of item groups explains the low entropy of companies into a few categories. High entropy indicates the participation of some companies in various tenders always without success.}

\section{Another dimension of the analysis is related to identifying PBUs involved in tenders with the participation of frequent losers. PBUs involved in ‘frequent losers tenders’ are also concentrated, as shown in Figure 8.}

\vspace{2\baselineskip}
\begin{figure}[H]
\centering
\includegraphics[width=12.73cm,height=9.55cm]{./images/image8.png}
\caption{Frequent Losers and PBUs}
\label{fig:frequent_losers_and_pbus}
\end{figure}


\vspace{1\baselineskip}
\ \ \ \ 

\section{Frequent Losers and Outcomes in Tenders}

\section{The outcomes observed in public tenders with participation by frequent losers can identify evidence of explicit collusion, potentially serving as a screening mechanism for the existence of these practices. As developed in Section 2, an essential principle of effective screening models is their ability to distinguish collusive behaviors from those that could arise from different competitive models.}

\section{This section investigates the properties of this collusion marker and observes the relationship between the presence of frequent losers and the performance of bidding contests in several dimensions, such as prices, the number of participants, and the number of bids. Based on these results, we evaluate the effectiveness of this marker in identifying collusive behavior.}

\section{Possible differences in bidding results are estimated by comparing public tenders with at least one frequent loser and tenders without the presence of a frequent loser. This estimation of differences in outcomes y between purchases i, with and without the participation of a frequent loser, of an item g and in year t has the following baseline specification:}

\vspace{1\baselineskip}
\begin{equation}
\mathbf{y}_{\mathbf{igt}}\mathbf{ = }\mathbf{\beta }\mathbf{losers}_{\mathbf{igt}}\mathbf{+}\mathbf{\alpha }_{\mathbf{g}}\mathbf{+}\mathbf{\lambda }_{\mathbf{t}}\mathbf{+}\mathbf{x\delta }\mathbf{+}\mathbf{\epsilon }_{\mathbf{igt,}}\mathbf{     }\left(\mathbf{2}\right)
\end{equation}


\vspace{1\baselineskip}
\section{where and are fixed effects of purchased items and year dummies, respectively. The variable has a value equal to 1 if the tender has at least one frequent loser and 0 otherwise. Additionally, are the control variables, such as dummies for purchasing units and tender mode.}

\section{The data used come from public purchases made at BEC from January 2009 to December 2019. However, for the purposes of this section, the subsample used in the estimates consists only of items with at least one purchasing process with the presence of frequent losers. The estimates for the negotiated prices are presented in Table 2.}

\begin{table}[H]
\begin{adjustbox}{max width=\textwidth}
\begin{tabular}{p{3.92cm}p{2.57cm}p{2.5cm}p{2.5cm}p{2.5cm}}
\hline
\multicolumn{1}{|p{5.3cm}}{} & 
\multicolumn{1}{|p{3.14cm}}{\centering
 (1)} & 
\multicolumn{1}{|p{3.02cm}}{\centering
 (2)} & 
\multicolumn{1}{|p{3.02cm}}{\centering
 (3)} & 
\multicolumn{1}{|p{3.02cm}|}{\centering
 (4)} \\ 
\hline
\multicolumn{1}{|p{5.3cm}}{} & 
\multicolumn{1}{|p{3.14cm}}{\centering
    General} & 
\multicolumn{1}{|p{3.02cm}}{\centering
    General} & 
\multicolumn{1}{|p{3.02cm}}{\centering
 Pregao} & 
\multicolumn{1}{|p{3.02cm}|}{\centering
    Convite} \\ 
\hline
\multicolumn{1}{|p{5.3cm}}{ losers} & 
\multicolumn{1}{|p{3.14cm}}{\centering
.1255$\ast$$\ast$$\ast$} & 
\multicolumn{1}{|p{3.02cm}}{\centering
.1011$\ast$$\ast$$\ast$} & 
\multicolumn{1}{|p{3.02cm}}{\centering
.1344$\ast$$\ast$$\ast$} & 
\multicolumn{1}{|p{3.02cm}|}{\centering
.0315$\ast$$\ast$$\ast$} \\ 
\hline
\multicolumn{1}{|p{5.3cm}}{} & 
\multicolumn{1}{|p{3.14cm}}{\centering
(.014)} & 
\multicolumn{1}{|p{3.02cm}}{\centering
(.0122)} & 
\multicolumn{1}{|p{3.02cm}}{\centering
(.0184)} & 
\multicolumn{1}{|p{3.02cm}|}{\centering
(.0046)} \\ 
\hline
\multicolumn{1}{|p{5.3cm}}{ convite} & 
\multicolumn{1}{|p{3.14cm}}{\centering
-.0449$\ast$$\ast$$\ast$} & 
\multicolumn{1}{|p{3.02cm}}{\centering
-.0442$\ast$$\ast$$\ast$} & 
\multicolumn{1}{|p{3.02cm}}{} & 
\multicolumn{1}{|p{3.02cm}|}{} \\ 
\hline
\multicolumn{1}{|p{5.3cm}}{} & 
\multicolumn{1}{|p{3.14cm}}{\centering
(.0057)} & 
\multicolumn{1}{|p{3.02cm}}{\centering
(.0056)} & 
\multicolumn{1}{|p{3.02cm}}{} & 
\multicolumn{1}{|p{3.02cm}|}{} \\ 
\hline
\multicolumn{1}{|p{5.3cm}}{ pregao} & 
\multicolumn{1}{|p{3.14cm}}{\centering
.4865$\ast$$\ast$$\ast$} & 
\multicolumn{1}{|p{3.02cm}}{\centering
.5653$\ast$$\ast$$\ast$} & 
\multicolumn{1}{|p{3.02cm}}{} & 
\multicolumn{1}{|p{3.02cm}|}{} \\ 
\hline
\multicolumn{1}{|p{5.3cm}}{} & 
\multicolumn{1}{|p{3.14cm}}{\centering
(.0177)} & 
\multicolumn{1}{|p{3.02cm}}{\centering
(.0208)} & 
\multicolumn{1}{|p{3.02cm}}{} & 
\multicolumn{1}{|p{3.02cm}|}{} \\ 
\hline
\multicolumn{1}{|p{5.3cm}}{ lquantity} & 
\multicolumn{1}{|p{3.14cm}}{\centering
-.2666$\ast$$\ast$$\ast$} & 
\multicolumn{1}{|p{3.02cm}}{\centering
-.2988$\ast$$\ast$$\ast$} & 
\multicolumn{1}{|p{3.02cm}}{\centering
-.4781$\ast$$\ast$$\ast$} & 
\multicolumn{1}{|p{3.02cm}|}{\centering
-.2118$\ast$$\ast$$\ast$} \\ 
\hline
\multicolumn{1}{|p{5.3cm}}{} & 
\multicolumn{1}{|p{3.14cm}}{\centering
(.0098)} & 
\multicolumn{1}{|p{3.02cm}}{\centering
(.0098)} & 
\multicolumn{1}{|p{3.02cm}}{\centering
(.0175)} & 
\multicolumn{1}{|p{3.02cm}|}{\centering
(.0109)} \\ 
\hline
\multicolumn{1}{|p{5.3cm}}{ \_cons} & 
\multicolumn{1}{|p{3.14cm}}{\centering
3.6905$\ast$$\ast$$\ast$} & 
\multicolumn{1}{|p{3.02cm}}{\centering
4.0337$\ast$$\ast$$\ast$} & 
\multicolumn{1}{|p{3.02cm}}{\centering
5.5803$\ast$$\ast$$\ast$} & 
\multicolumn{1}{|p{3.02cm}|}{\centering
3.7513$\ast$$\ast$$\ast$} \\ 
\hline
\multicolumn{1}{|p{5.3cm}}{} & 
\multicolumn{1}{|p{3.14cm}}{\centering
(.0352)} & 
\multicolumn{1}{|p{3.02cm}}{\centering
(.0703)} & 
\multicolumn{1}{|p{3.02cm}}{\centering
(.1422)} & 
\multicolumn{1}{|p{3.02cm}|}{\centering
(.0825)} \\ 
\hline
\multicolumn{1}{|p{5.3cm}}{ Observations} & 
\multicolumn{1}{|p{3.14cm}}{\centering
1671773} & 
\multicolumn{1}{|p{3.02cm}}{\centering
1671773} & 
\multicolumn{1}{|p{3.02cm}}{\centering
474219} & 
\multicolumn{1}{|p{3.02cm}|}{\centering
924725} \\ 
\hline
\multicolumn{1}{|p{5.3cm}}{ R-squared} & 
\multicolumn{1}{|p{3.14cm}}{\centering
.2085} & 
\multicolumn{1}{|p{3.02cm}}{\centering
.2596} & 
\multicolumn{1}{|p{3.02cm}}{\centering
.4219} & 
\multicolumn{1}{|p{3.02cm}|}{\centering
.2066} \\ 
\hline
\multicolumn{1}{|p{5.3cm}}{Item Dummies} & 
\multicolumn{1}{|p{3.14cm}}{\centering
YES} & 
\multicolumn{1}{|p{3.02cm}}{\centering
YES} & 
\multicolumn{1}{|p{3.02cm}}{\centering
YES} & 
\multicolumn{1}{|p{3.02cm}|}{\centering
YES} \\ 
\hline
\multicolumn{1}{|p{5.3cm}}{Year Dummies} & 
\multicolumn{1}{|p{3.14cm}}{\centering
YES} & 
\multicolumn{1}{|p{3.02cm}}{\centering
YES} & 
\multicolumn{1}{|p{3.02cm}}{\centering
YES} & 
\multicolumn{1}{|p{3.02cm}|}{\centering
YES} \\ 
\hline
\multicolumn{1}{|p{5.3cm}}{PBU Dummies} & 
\multicolumn{1}{|p{3.14cm}}{\centering
NO} & 
\multicolumn{1}{|p{3.02cm}}{\centering
YES} & 
\multicolumn{1}{|p{3.02cm}}{\centering
YES} & 
\multicolumn{1}{|p{3.02cm}|}{\centering
YES} \\ 
\hline
\multicolumn{5}{|p{17.49cm}|}{\textit{Standard errors are in parentheses}} \\ 
\hline
\multicolumn{5}{|p{17.49cm}|}{\textit{$\ast$$\ast$$\ast$ p<.01, $\ast$$\ast$ p<.05, $\ast$ p<.1}} \\ 
\hline
\end{tabular}
\end{adjustbox}
\caption{With vs. Without Frequent Losers}
\label{tab:vs_frequent_losers}\end{table}
\vspace{17\baselineskip}
\section{It can be observed that the negotiated prices are consistently higher in tenders with the presence of frequent losers. Considering variations in the baseline specification containing all the bidding modes, the negotiated prices, on average, are between 10.64% and 13.37% higher in bids with losers.}

\section{The results also suggest that prices in pregoes tend to be more affected than prices in convites by the participation of frequent losers. Considering only pregoes, it is possible to observe that prices tend to be 14.39% higher in bids that include losers than when losers are not present; for convites, prices are only 3.20% higher on the same comparative basis.}

\section{Possible explanations for this difference might be related to (i) the degree of discretion of the public official in charge of the tender and (ii) the possibility for suppliers to affect the dynamics and outcome of the process.}

\section{In convites, the process consists of merely opening (virtual) envelopes containing participating suppliers’ proposals. After making the bids public, the auctioneer’s role is restricted to declaring the firm that offered the lowest price the winner. Thus, the interaction between the participating firms and public officers is minimal during the purchase procedure.}

\section{In the case of the pregao, the public officer in charge has a more active role in the process, and there is more room for firms to influence the process. In addition to revealing the initial bids in the first phase, the auctioneer is responsible for coordinating and monitoring firms’ bids during the auction phase.}

\section{During this auction phase, firms and PBUs interact through the real-time bidding process and real-time chat. A great variety of information is exchanged through this communication channel, such as confirming product specifications, complaints about competitors’ performance, and even dissatisfaction with the way the contest is conducted.}

\section{At the end of the auction phase, there is another moment of interaction between suppliers and PBUs. There is an ex-post negotiation phase in which public and private parties can direct bargaining for the lowest price. Thus, in pregao, there is more room for manipulation. This situation may lead to higher prices in this type of tender.}

\section{Higher prices in tenders with frequent losers may suggest that there is less competitiveness or aggressiveness for suppliers. This fact cannot be verified in the estimates of Tables 3 and 4. In these tables, we compare tenders with and without frequent losers in terms of the number of firms and bids, respectively.}

\vspace{1\baselineskip}
\begin{center}
Table 3. Number of Firms: With vs. Without Frequent Losers
\end{center}


\begin{table}[H]
\begin{adjustbox}{max width=\textwidth}
\begin{tabular}{p{3.98cm}p{2.52cm}p{2.53cm}p{2.53cm}p{2.45cm}}
\hline
\multicolumn{1}{|p{5.3cm}}{} & 
\multicolumn{1}{|p{3.14cm}}{\centering
 (1)} & 
\multicolumn{1}{|p{3.02cm}}{\centering
 (2)} & 
\multicolumn{1}{|p{3.02cm}}{\centering
 (3)} & 
\multicolumn{1}{|p{3.02cm}|}{\centering
 (4)} \\ 
\hline
\multicolumn{1}{|p{5.3cm}}{} & 
\multicolumn{1}{|p{3.14cm}}{\centering
    General} & 
\multicolumn{1}{|p{3.02cm}}{\centering
    General} & 
\multicolumn{1}{|p{3.02cm}}{\centering
 Pregao} & 
\multicolumn{1}{|p{3.02cm}|}{\centering
    Convite} \\ 
\hline
\multicolumn{1}{|p{5.3cm}}{ losers} & 
\multicolumn{1}{|p{3.14cm}}{\centering
.329$\ast$$\ast$$\ast$} & 
\multicolumn{1}{|p{3.02cm}}{\centering
.3256$\ast$$\ast$$\ast$} & 
\multicolumn{1}{|p{3.02cm}}{\centering
.3104$\ast$$\ast$$\ast$} & 
\multicolumn{1}{|p{3.02cm}|}{\centering
.2926$\ast$$\ast$$\ast$} \\ 
\hline
\multicolumn{1}{|p{5.3cm}}{} & 
\multicolumn{1}{|p{3.14cm}}{\centering
(.0052)} & 
\multicolumn{1}{|p{3.02cm}}{\centering
(.0048)} & 
\multicolumn{1}{|p{3.02cm}}{\centering
(.008)} & 
\multicolumn{1}{|p{3.02cm}|}{\centering
(.0041)} \\ 
\hline
\multicolumn{1}{|p{5.3cm}}{ convite} & 
\multicolumn{1}{|p{3.14cm}}{\centering
-.7742$\ast$$\ast$$\ast$} & 
\multicolumn{1}{|p{3.02cm}}{\centering
-.7755$\ast$$\ast$$\ast$} & 
\multicolumn{1}{|p{3.02cm}}{} & 
\multicolumn{1}{|p{3.02cm}|}{} \\ 
\hline
\multicolumn{1}{|p{5.3cm}}{} & 
\multicolumn{1}{|p{3.14cm}}{\centering
(.0069)} & 
\multicolumn{1}{|p{3.02cm}}{\centering
(.007)} & 
\multicolumn{1}{|p{3.02cm}}{} & 
\multicolumn{1}{|p{3.02cm}|}{} \\ 
\hline
\multicolumn{1}{|p{5.3cm}}{ pregao} & 
\multicolumn{1}{|p{3.14cm}}{\centering
-.1181$\ast$$\ast$$\ast$} & 
\multicolumn{1}{|p{3.02cm}}{\centering
-.1228$\ast$$\ast$$\ast$} & 
\multicolumn{1}{|p{3.02cm}}{} & 
\multicolumn{1}{|p{3.02cm}|}{} \\ 
\hline
\multicolumn{1}{|p{5.3cm}}{} & 
\multicolumn{1}{|p{3.14cm}}{\centering
(.0081)} & 
\multicolumn{1}{|p{3.02cm}}{\centering
(.0099)} & 
\multicolumn{1}{|p{3.02cm}}{} & 
\multicolumn{1}{|p{3.02cm}|}{} \\ 
\hline
\multicolumn{1}{|p{5.3cm}}{ lquantity} & 
\multicolumn{1}{|p{3.14cm}}{\centering
.1504$\ast$$\ast$$\ast$} & 
\multicolumn{1}{|p{3.02cm}}{\centering
.151$\ast$$\ast$$\ast$} & 
\multicolumn{1}{|p{3.02cm}}{\centering
.0925$\ast$$\ast$$\ast$} & 
\multicolumn{1}{|p{3.02cm}|}{\centering
.1817$\ast$$\ast$$\ast$} \\ 
\hline
\multicolumn{1}{|p{5.3cm}}{} & 
\multicolumn{1}{|p{3.14cm}}{\centering
(.0021)} & 
\multicolumn{1}{|p{3.02cm}}{\centering
(.0022)} & 
\multicolumn{1}{|p{3.02cm}}{\centering
(.0028)} & 
\multicolumn{1}{|p{3.02cm}|}{\centering
(.0028)} \\ 
\hline
\multicolumn{1}{|p{5.3cm}}{ \_cons} & 
\multicolumn{1}{|p{3.14cm}}{\centering
.9159$\ast$$\ast$$\ast$} & 
\multicolumn{1}{|p{3.02cm}}{\centering
1.2013$\ast$$\ast$$\ast$} & 
\multicolumn{1}{|p{3.02cm}}{\centering
1.0851$\ast$$\ast$$\ast$} & 
\multicolumn{1}{|p{3.02cm}|}{\centering
.8704$\ast$$\ast$$\ast$} \\ 
\hline
\multicolumn{1}{|p{5.3cm}}{} & 
\multicolumn{1}{|p{3.14cm}}{\centering
(.0093)} & 
\multicolumn{1}{|p{3.02cm}}{\centering
(.028)} & 
\multicolumn{1}{|p{3.02cm}}{\centering
(.0451)} & 
\multicolumn{1}{|p{3.02cm}|}{\centering
(.0624)} \\ 
\hline
\multicolumn{1}{|p{5.3cm}}{ Observations} & 
\multicolumn{1}{|p{3.14cm}}{\centering
1670719} & 
\multicolumn{1}{|p{3.02cm}}{\centering
1670719} & 
\multicolumn{1}{|p{3.02cm}}{\centering
473819} & 
\multicolumn{1}{|p{3.02cm}|}{\centering
924336} \\ 
\hline
\multicolumn{1}{|p{5.3cm}}{ R-squared} & 
\multicolumn{1}{|p{3.14cm}}{\centering
.3675} & 
\multicolumn{1}{|p{3.02cm}}{\centering
.3882} & 
\multicolumn{1}{|p{3.02cm}}{\centering
.2176} & 
\multicolumn{1}{|p{3.02cm}|}{\centering
.2699} \\ 
\hline
\multicolumn{1}{|p{5.3cm}}{Items Dummies} & 
\multicolumn{1}{|p{3.14cm}}{\centering
YES} & 
\multicolumn{1}{|p{3.02cm}}{\centering
YES} & 
\multicolumn{1}{|p{3.02cm}}{\centering
YES} & 
\multicolumn{1}{|p{3.02cm}|}{\centering
YES} \\ 
\hline
\multicolumn{1}{|p{5.3cm}}{Year Dummies} & 
\multicolumn{1}{|p{3.14cm}}{\centering
YES} & 
\multicolumn{1}{|p{3.02cm}}{\centering
YES} & 
\multicolumn{1}{|p{3.02cm}}{\centering
YES} & 
\multicolumn{1}{|p{3.02cm}|}{\centering
YES} \\ 
\hline
\multicolumn{1}{|p{5.3cm}}{PBU Dummies} & 
\multicolumn{1}{|p{3.14cm}}{\centering
NO} & 
\multicolumn{1}{|p{3.02cm}}{\centering
YES} & 
\multicolumn{1}{|p{3.02cm}}{\centering
YES} & 
\multicolumn{1}{|p{3.02cm}|}{\centering
YES} \\ 
\hline
\multicolumn{5}{|p{17.49cm}|}{\textit{Standard errors are in parentheses}} \\ 
\hline
\multicolumn{5}{|p{17.49cm}|}{\textit{$\ast$$\ast$$\ast$ p<.01, $\ast$$\ast$ p<.05, $\ast$ p<.1}} \\ 
\hline
\end{tabular}
\end{adjustbox}
\end{table}
\vspace{18\baselineskip}
\section{More firms participate and more bids are offered in tenders with frequent losers than in bids without their presence. Between 33.99% and 38.96% more firms participate in tenders with losers, with an average of 25.51% to 34.35% more bids offered.}

\vspace{2\baselineskip}
\begin{center}
Table 4. Number of Bids: With vs. Without Frequent Losers
\end{center}


\begin{table}[H]
\begin{adjustbox}{max width=\textwidth}
\begin{tabular}{p{4.21cm}p{2.5cm}p{2.4cm}p{2.4cm}p{2.5cm}}
\hline
\multicolumn{1}{|p{4.21cm}}{} & 
\multicolumn{1}{|p{2.5cm}}{\centering
 (1)} & 
\multicolumn{1}{|p{2.4cm}}{\centering
 (2)} & 
\multicolumn{1}{|p{2.4cm}}{\centering
 (3)} & 
\multicolumn{1}{|p{2.5cm}|}{\centering
 (4)} \\ 
\hline
\multicolumn{1}{|p{4.21cm}}{} & 
\multicolumn{1}{|p{2.5cm}}{\centering
    General} & 
\multicolumn{1}{|p{2.4cm}}{\centering
    General} & 
\multicolumn{1}{|p{2.4cm}}{\centering
 Pregao} & 
\multicolumn{1}{|p{2.5cm}|}{\centering
    Convite} \\ 
\hline
\multicolumn{1}{|p{4.21cm}}{ losers} & 
\multicolumn{1}{|p{2.5cm}}{\centering
.2953$\ast$$\ast$$\ast$} & 
\multicolumn{1}{|p{2.4cm}}{\centering
.2937$\ast$$\ast$$\ast$} & 
\multicolumn{1}{|p{2.4cm}}{\centering
.2272$\ast$$\ast$$\ast$} & 
\multicolumn{1}{|p{2.5cm}|}{\centering
.2926$\ast$$\ast$$\ast$} \\ 
\hline
\multicolumn{1}{|p{4.21cm}}{} & 
\multicolumn{1}{|p{2.5cm}}{\centering
(.0062)} & 
\multicolumn{1}{|p{2.4cm}}{\centering
(.006)} & 
\multicolumn{1}{|p{2.4cm}}{\centering
(.0108)} & 
\multicolumn{1}{|p{2.5cm}|}{\centering
(.0041)} \\ 
\hline
\multicolumn{1}{|p{4.21cm}}{ convite} & 
\multicolumn{1}{|p{2.5cm}}{\centering
.0314$\ast$$\ast$$\ast$} & 
\multicolumn{1}{|p{2.4cm}}{\centering
.0249$\ast$$\ast$$\ast$} & 
\multicolumn{1}{|p{2.4cm}}{} & 
\multicolumn{1}{|p{2.5cm}|}{} \\ 
\hline
\multicolumn{1}{|p{4.21cm}}{} & 
\multicolumn{1}{|p{2.5cm}}{\centering
(.0093)} & 
\multicolumn{1}{|p{2.4cm}}{\centering
(.0093)} & 
\multicolumn{1}{|p{2.4cm}}{} & 
\multicolumn{1}{|p{2.5cm}|}{} \\ 
\hline
\multicolumn{1}{|p{4.21cm}}{ pregao} & 
\multicolumn{1}{|p{2.5cm}}{\centering
.931$\ast$$\ast$$\ast$} & 
\multicolumn{1}{|p{2.4cm}}{\centering
.9225$\ast$$\ast$$\ast$} & 
\multicolumn{1}{|p{2.4cm}}{} & 
\multicolumn{1}{|p{2.5cm}|}{} \\ 
\hline
\multicolumn{1}{|p{4.21cm}}{} & 
\multicolumn{1}{|p{2.5cm}}{\centering
(.01)} & 
\multicolumn{1}{|p{2.4cm}}{\centering
(.0126)} & 
\multicolumn{1}{|p{2.4cm}}{} & 
\multicolumn{1}{|p{2.5cm}|}{} \\ 
\hline
\multicolumn{1}{|p{4.21cm}}{ lquantity} & 
\multicolumn{1}{|p{2.5cm}}{\centering
.169$\ast$$\ast$$\ast$} & 
\multicolumn{1}{|p{2.4cm}}{\centering
.1674$\ast$$\ast$$\ast$} & 
\multicolumn{1}{|p{2.4cm}}{\centering
.1059$\ast$$\ast$$\ast$} & 
\multicolumn{1}{|p{2.5cm}|}{\centering
.1817$\ast$$\ast$$\ast$} \\ 
\hline
\multicolumn{1}{|p{4.21cm}}{} & 
\multicolumn{1}{|p{2.5cm}}{\centering
(.0026)} & 
\multicolumn{1}{|p{2.4cm}}{\centering
(.0027)} & 
\multicolumn{1}{|p{2.4cm}}{\centering
(.004)} & 
\multicolumn{1}{|p{2.5cm}|}{\centering
(.0028)} \\ 
\hline
\multicolumn{1}{|p{4.21cm}}{ \_cons} & 
\multicolumn{1}{|p{2.5cm}}{\centering
.7849$\ast$$\ast$$\ast$} & 
\multicolumn{1}{|p{2.4cm}}{\centering
1.5003$\ast$$\ast$$\ast$} & 
\multicolumn{1}{|p{2.4cm}}{\centering
2.17$\ast$$\ast$$\ast$} & 
\multicolumn{1}{|p{2.5cm}|}{\centering
.8704$\ast$$\ast$$\ast$} \\ 
\hline
\multicolumn{1}{|p{4.21cm}}{} & 
\multicolumn{1}{|p{2.5cm}}{\centering
(.0114)} & 
\multicolumn{1}{|p{2.4cm}}{\centering
(.0478)} & 
\multicolumn{1}{|p{2.4cm}}{\centering
(.077)} & 
\multicolumn{1}{|p{2.5cm}|}{\centering
(.0624)} \\ 
\hline
\multicolumn{1}{|p{4.21cm}}{ Observations} & 
\multicolumn{1}{|p{2.5cm}}{\centering
1670719} & 
\multicolumn{1}{|p{2.4cm}}{\centering
1670719} & 
\multicolumn{1}{|p{2.4cm}}{\centering
473819} & 
\multicolumn{1}{|p{2.5cm}|}{\centering
924336} \\ 
\hline
\multicolumn{1}{|p{4.21cm}}{ R-squared} & 
\multicolumn{1}{|p{2.5cm}}{\centering
.2945} & 
\multicolumn{1}{|p{2.4cm}}{\centering
.3117} & 
\multicolumn{1}{|p{2.4cm}}{\centering
.1351} & 
\multicolumn{1}{|p{2.5cm}|}{\centering
.2698} \\ 
\hline
\multicolumn{1}{|p{4.21cm}}{Items Dummies} & 
\multicolumn{1}{|p{2.5cm}}{\centering
YES} & 
\multicolumn{1}{|p{2.4cm}}{\centering
YES} & 
\multicolumn{1}{|p{2.4cm}}{\centering
YES} & 
\multicolumn{1}{|p{2.5cm}|}{\centering
YES} \\ 
\hline
\multicolumn{1}{|p{4.21cm}}{Year Dummies} & 
\multicolumn{1}{|p{2.5cm}}{\centering
YES} & 
\multicolumn{1}{|p{2.4cm}}{\centering
YES} & 
\multicolumn{1}{|p{2.4cm}}{\centering
YES} & 
\multicolumn{1}{|p{2.5cm}|}{\centering
YES} \\ 
\hline
\multicolumn{1}{|p{4.21cm}}{PBU Dummies} & 
\multicolumn{1}{|p{2.5cm}}{\centering
NO} & 
\multicolumn{1}{|p{2.4cm}}{\centering
YES} & 
\multicolumn{1}{|p{2.4cm}}{\centering
YES} & 
\multicolumn{1}{|p{2.5cm}|}{\centering
YES} \\ 
\hline
\multicolumn{5}{|p{14.01cm}|}{\textit{Standard errors are in parentheses}} \\ 
\hline
\multicolumn{5}{|p{14.01cm}|}{\textit{$\ast$$\ast$$\ast$ p<.01, $\ast$$\ast$ p<.05, $\ast$ p<.1}} \\ 
\hline
\end{tabular}
\end{adjustbox}
\end{table}
\vspace{19\baselineskip}
\section{These results are compatible with the literature on cartels and other illicit purchases. As ) notes, cartel participants’ behavior also responds strategically to competition policy enforcement. A higher number of participating firms and bidders could be associated with attempting to build a mechanism to conceal possible illicit competitive practices.}

\section{It is important to note that, assuming there is no correlation between the companies’ cost structure and the presence of frequent losers, the simultaneous occurrence of higher prices and a higher number of bids and participants in public tenders is inconsistent with expectations for a model of competition or tacit collusion.}

\section{This scenario might mean that the screening proposal based on frequent losers presents two essential properties. First, it shows that it is capable of distinguishing collusive behaviors from those that would be expected in competition. Second and more importantly, because the observed behavior reveals the intention to hide the collusion, it is a marker that identifies explicit and non-tacit collusion, a recurrent problem in screening models ).}

\section{To establish the reliability of this screening model, however, it is necessary to highlight possible alternative explanations for the presence of frequent losers in public procurement. A first alternative hypothesis is the existence of a learning curve for companies to become victorious at BEC. Supplying goods and services to the state would require suppliers to develop specific processes and legal adjustments that would demand consistent efforts over some time.}

\section{If there were such a learning curve, then companies that eventually win the first bid would spend a considerable period of time losing until they become victorious. This situation may be the case for some suppliers; however, the data suggest that this is not the general case. Figure 9 shows the number of contests up to a company’s first victory, limited to 100 contests for better viewing. Almost 30% of companies win one of the first three events in which they participate.}

\vspace{1\baselineskip}
\begin{figure}[H]
\includegraphics[width=10.82cm,height=7.3cm]{./images/image9.png}
\caption{-Number of Participations in Tenders until the First Victory}
\label{fig:number_participations_tenders_first_victory}
\end{figure}


\vspace{1\baselineskip}
\section{In addition, as explained in Section 3, it is known that a small number of companies win public tenders. Thus, a learning curve hypothesis seems unlikely, except for a few companies that achieve very frequent wins in bids.}

\section{There is an additional alternative explanation that could be considered in this context. Frequent losers might be a phenomenon found mostly in the short term: there would be a significant number of participations without victories in a short period, followed by abandonment of the public contracting market. As shown in Figure 3, presented in Section 3, companies have participated in tenders for up to ten years and were repeatedly defeated. It is not reasonable to suppose that there is a financial incentive to participate in tenders in good faith and to lose repeatedly.}

\section{It is also important to note that even if the companies participated for a short time, they still participated in a sufficient number of tenders to be classified as frequent losers by the outlier detection method.}

\vspace{2\baselineskip}
\section{Conclusion}

\section{Companies that systematically lose public tenders in which they participate are not expected in a competitive environment. An analysis of the relationship between performance variables and the participation of these companies in public bids in São Paulo State between 2009 and 2019 indicates that the negotiated price was 10% higher; however, the number of participants was 32% higher, and the number of bids was 29% higher.}

\section{This apparent paradox, in which one variable indicates less competitive pressure (i.e., negotiated prices) while others indicate more competition (i.e., number of participants and number of bids), is consistent with the behavior of a cartel that seeks to avoid detection. This scenario is the function that would intuitively be expected of a frequent loser, that is, to simulate a higher level of competition, drive away competitors, and avoid detection by control agencies. Precisely because of this characteristic, frequent losers can be used as markers in a cartel screening method.}

\section{This proposal presents some virtues common to other screening methods and addresses two limitations common to most available models. It is a simple application method that requires only data already available to competition authorities and control agencies.}

\section{The identification of frequent losers may occur before the public tender on the occasion of the first stage of the bid, in which the participants are defined. Unlike all available screening methods based on prices and quantities observed after the cartel’s materialization, frequent losers may signal suspicious behavior before bidding occurs, reducing the costs associated with a defrauded event.}

\section{This property stems from a characteristic of cartels in public tenders: the agreement among bidders occurs before the tender takes place, and it is possible to observe elements that may signal fraud before the conclusion of the bid. The presence of frequent losers also emphasizes the separation between explicit collusion and tacit collusion since the behavior identified is consistent with the deliberate act of concealing competition.}

\section{There are, however, some limitations in the use of these screening methods; some of these limitations are remedial, and another is common to all. The proposed method identifies suspicious conduct only in cases of bidding cartels, where some mechanism for transferring the benefits of the cartel among its members is possible. This is still a relevant subset of cases, but it does not apply to all cartels. However, this limitation is remedial because the method can be associated with any other screening method and can act in a complementary manner, as suggested by .}

\section{common problem with all screening methods is that once they become known to cartel participants, the cartel participants will modify their behavior to avoid detection (Schinkel, 2014). However, this is not a reason to rule out screening methods; on the contrary, it is necessary to develop them continuously, and the cumulative process of screening methods is associated with increasing costs for cartel participants to avoid detection. Eventually, this cumulative process may result in the deterrence of such an antitrust offense.}


\newpage

\printbibliography

\end{document}